%%%%%%%%%%%%%%%%%%%%%%%%%%%%%%%%%%%%%%%%%%%%%%%%%%%%%
%			HLAVIČKA								%
%%%%%%%%%%%%%%%%%%%%%%%%%%%%%%%%%%%%%%%%%%%%%%%%%%%%%
\documentclass[openany,12pt]{memoir}

\usepackage[utf8]{inputenc} 
\usepackage[czech]{babel}
\usepackage[T1]{fontenc}
\usepackage[top=1.5cm, bottom=2cm, left=2cm, right=2cm]{geometry}  % --> NASTAVENÍ OKRAJŮ
\usepackage{fancyhdr}
\usepackage{graphicx}
%\usepackage{lmodern}
\usepackage{tgschola}  % --- Silnější font
\usepackage{xwatermark}
\usepackage{xcolor}
\usepackage{changepage}
\usepackage{newunicodechar}



%%%%%% Package na zpěvník
\usepackage[full]{leadsheets}%http://mirrors.nic.cz/tex-archive/macros/latex/contrib/leadsheets/leadsheets_en.pdf   --> dokumentace	
\definesongtitletemplate{empty}{} 
\setchords{
format = \bfseries,   %tučné akordy
minor = {mi},% 
input-notation = {german},%
output-notation = {german}%
}
\definesongtitletemplate{empty}{} 

\newlength{\drop}
\newwatermark[allpages,color=red!50,angle=0,scale=2, xpos=0,ypos=0]{\includegraphics[width=5cm]{obr/pozadi2.jpg}} %--> dvojka na pozadí


%%%% Vlastní příkazy
\newcounter{Slokočet}   %Automatické číslování slok
\newcommand{\mezera}{\vspace*{0.5cm}}   %Horizontální odsazení slok
\newcommand{\stred}{5.2cm}   %%% Na zarovnání slok doprostřed, pozn. automatičtější zarovnávání na střed nejde
\newcommand{\refren}{\mezera \noindent \textbf{R:} } %refrén
\newcommand{\sloka}{\mezera \noindent \addtocounter{Slokočet}{1} \arabic{Slokočet}. } 	%sloka, která se automaticky čísluje
\newcommand{\ssloka}{\mezera \noindent}  % vlastní číslo sloky
\newcommand{\carka}{,\:}
\newcommand{\m}[1]{\color{white}{#1}}  %Pro akordy
\newcommand{\ap}{'} %Pro apostrof


%%%%%%%%%%%%%%%%%%%%%%%%%%%%%%%%%%%%%%%%%%%%%%%%%%%%%%
%			NÁVOD									 %
%%%%%%%%%%%%%%%%%%%%%%%%%%%%%%%%%%%%%%%%%%%%%%%%%%%%%%
%1. Věci v hlavičce IGNOROVAT
%2. Píseň psát do prostoru mezi \begin{song} a \end{song}
%3. další řádek se značí dvěma odsazeníma (= dvakrát stisknout enter)
%4. \refren vždy na začátku refrenu a \sloka na začátek sloky (automaticky se čísluje)
%  \ = alt gr + q ; [/] = alt gr f/g ; {/} = alt gr + b/n; ^ = alt gr + 3 + mezera
%Cokoliv napíšete do ^{  } se bude brát jako akord
%když se toto bude dotýkat nějakého slova (nebude mezi tím a slovem mezera)
%tak se akordy zjeví nad slovem, ale pište to před slovo
%Když se to nedotýká slova, tak akord lítá ve vzduchu a vytiskne se větší mezera
%První možnost je asi preferovanější
%5. Akordy stačí psát jen do první sloky, když se nezmění -- kytaristi to zvládnou



%\usepackage{subfiles}
%</++>

\begin{document}

\pagestyle{empty}
%\addtocontents{toc}{\protect\thispagestyle{empty}}
%\tableofcontents \thispagestyle{empty}\newpage
\newgeometry{top=1.5cm, bottom = 0cm, left = 2cm, right = 2cm}

\newpage

%Pro kompilaci je důležitý, aby ten input a \end{document} zůstali na posledních
%dvou řádcích!!!!
\begin{song}{title=\predtitle\centering Chci ti říct \\\large Mig 21  \vspace*{-0.3cm}}  %% sem se napíše jméno songu a autor
\begin{centerjustified}


\predehra
4$\times $\textbf{Gmi C2}

\sloka
	^{Gmi}Chladná ^{C2}rána jsou ^{Gmi}tomu, kdo ^{C2}strádá,

	^{Gmi\,\,}láskou ^{C2\z }bolestnou ^{Gmi\z C2}neopětovanou.

	Pokud ^{Gmi}víš, jak ^{C2}rád tě mám, ^{Gmi}a ty mě máš ^{C2}ráda.

	^{Gmi}Chuť a ^{C2}vůni tvou ^{Gmi}nejde ^{C2\z }zapomenout.

\sloka
	Chystá ^{Dmi}se prej konec ^{C2}světa, ^{Gmi}tramtadá

	a ^{Dmi}mě jen jedna ^{C2}věta ^{\z B2}napadá. (\dots\,tak pojď s tim ven!)

\refren
	^{F}Chci ti říct, že mám tě ^{C2}rád, že ^{Gmi}miluju tě jako blázen,

	^{F}chci si nechat ^{C2\z }tetovat jen ^{Gmi}jméno tvé, chci vypít bazén

	^{F}vody, kterou ^{C2\z }plavala jsi, ^{Gmi\z }zasadit si tvoje vlasy,

	^{F\z }rozfoukat je ^{C2}do ulic ^{Gmi}chci ti říct, že rád tě mám ^{B7}víc.

\sloka
	Dýcháš na mou tvář, já na tvou dýchám.

	Chvěješ se myšlenkou, že se rozední.

	Děláš, jako že nespěcháš, já že nepospíchám.

	Nechce se z objetí, když je poslední.

\sloka  = 2.

\refren


\sloka
	/: ^{B7}Mít tak promyšlenej plán, předem fakt moc se omlouvám, 

	ale ode dneška od půl jedný musíme bejt zodpovědný. :/

	^{B7}Mít tak promyšlenej plán, ale nemám ho, nemám ho, nemám\elipsa.\elipsa.\elipsa.\,a tak:

\refren


\end{centerjustified}
\setcounter{Slokočet}{0}
\end{song}

\end{document}
